\section{Survival within the resource budgets}\label{sec:survival}

The preceding construction always respects its dictionary capacities,
but it may discard a bucket containing a solution or stop at the
operation cap. We now show that a fixed solution reaches a completed
check within the cap with inverse-polynomial probability, enough for
repetition to succeed. The order of the argument matters: first fix a
surviving representation, subtract an unconditional resource-failure
bound, and only then use the extra prime and separating families.

\subsection{Fixing a witness before the last random choices}

Throughout this section suppose that the input is a YES instance and
\(\mu<.00765n\), where \(\mu\) is the mean deficiency of a uniform
size-\(s\) subset of the input indices. We first use ideal exact sampling.
The extra prime \(q\) and the separating families are not exposed in this
subsection. Our objective is to select one representation before either
of those sources of randomness is used. We first balance the fixed
solution inside the mixers, retain many distinct choices with the
prescribed leaf counts, and finally show that the low residues hit
some of them. Every selection of a witness in this argument is an
analysis choice; the program enumerates candidates without knowing it.

\begin{lemma}\label{lem:representation}
For a main-branch YES input with \(\mu<.00765n\), ideal sampling
produces, with probability at least \(u^{-100}\), a passing spread test and a representation of a solution satisfying all prescribed leaf
counts and all low modular filters. Its two parent weights
\(x_L,x_R\), with \(x_L+x_R=t'\), and all of its masks can be fixed as
functions of the nonfamily choices preceding \(q\). Conditional on
these choices, \(q\) still has its prescribed independent prime law;
the separating families also retain their prescribed product law,
with the now fixed layout parameters.
\end{lemma}

\begin{proof}
\emph{Balance the solution and its deficiencies.}
Fix an original solution \(S\) for analysis. Under the random sign flip
\(F\), its signed counterpart \(S'=S\mathbin{\triangle}F\) is a uniform
subset of the input indices, independent of the mixers. Put
\(A_i=S'\cap M_i\). By Lemma~\ref{lem:entropy}, the balance event
\[
 |A_0|=|A_L|=|A_R|=s
\]
has probability at least \((m+1)^{-3}\). This event has the same
probability for every choice of the three mixers before ordering them.
Conditional on it, their distribution is unchanged, and their selected
halves are independent uniform halves given the mixers. Each selected
half and its complement in a pre-order mixer is therefore marginally a
uniform size-\(s\) set of input indices.

By Lemma~\ref{lem:deficiency}, with conditional probability
\(1-O(1/n)\), all six halves have deficiency at most
\(\tau=.0078n\): a deviation of \(.0001n\) above their common mean
already suffices. The same lemma, now for size-\(m\) sets, shows that
each pre-order mixer's deficiency is within \(.0002n\) of the
size-\(m\) mean, again with conditional failure probability \(O(1/n)\).
Thus the spread test \eqref{eq:spread} passes jointly with these events
for sufficiently large \(n\). These are statements about the original
weights, equally valid after sign changes by deficiency invariance.
For this analysis we may sample the signs even on choices on which the
algorithm would have stopped at the spread test.

If \(A,B\) partition a mixer \(M\) into halves, then
\(N(M)\le N(A)N(B)\), so
\[
 \defect(A)+\defect(B)\le \defect(M).
\]
At least one half has deficiency at most \(\defect(M)/2\). Given the
three unordered pairs of halves, each orientation as the selected half
is uniform, independently across the mixers. Both preceding good events concern the original-weight deficiencies
of the unordered halves and are invariant under these orientations,
as is the ordering of the mixers. With further conditional probability at least \(1/8\), then,
\begin{equation}\label{surv:small-deficiency}
 \defect(A_i)\le f_i=\min(\tau,d_i/2)
 \qquad(i=0,L,R).
\end{equation}
No conditional independence of the signs is asserted or needed here.
Sign invariance and \eqref{surv:small-deficiency} give at least
\(2^{s-f_i}=2^{L_i}\) distinct signed subset sums on \(A_i\).

\emph{Retain representatives with the prescribed leaf counts.}
Partition these sums according to subset size. Some size has at least
\(2^{L_i}/(s+1)\) distinct sums. Choose a maximizing size at most
\(s/2=m/4\), using complementation in \(A_i\), which preserves the
number of distinct sums between complementary sizes. This size exceeds
\(m/6\) for sufficiently large \(n\). Otherwise its number of subsets
would be at most \(2^{s h(1/3)}\), whereas
\[
 \frac{s-\tau-\log(s+1)}{s}\longrightarrow .935
 >\frac{14}{15}>h(1/3).
\]
For each mixer fix such a size and one subset for each distinct sum of
that size. These are analysis choices made before the size guesses and
subblock partitions. The probability that all three guesses \(k_i\)
equal these sizes is at least \((m+1)^{-3}\).

We next retain many of these representatives with the required finer
counts. Fix one representative \(D\subset A_i\), and write \(l=4\)
for the top mixer and \(l=2\) for either other mixer. In a uniform
partition into \(l\) labeled equal subblocks, prescribe the counts of
the three categories
\[
 D,\qquad A_i\setminus D,\qquad M_i\setminus A_i
\]
to be respectively the smaller-mask counts, their prescribed
complementary counts, and an equal split of the outside category.
These vectors fill every block exactly. In a top block their counts
are \(c_a,g/2-c_a,g/2\); in a local block they are
\(a_{jr},s/2-a_{jr},s/2\). The divisibility of \(m\) by eight makes
all these complementary counts integral. Each category's counts are
as equal as possible. For \(v\) objects,
a most balanced multinomial coefficient is maximal, since transferring
one object from a larger count to a count smaller by at least two
increases the coefficient. There are at most \((v+1)^l\) possible
count vectors, and their coefficients sum to \(l^v\). Each of our
three favorable coefficients is therefore at least
\(l^v/(v+1)^l\). Dividing their product by the multinomial coefficient
for the whole partition, which is at most \(l^m\), shows that \(D\)
survives with probability at least
\[
 (m+1)^{-3l}\ge u^{-3l}.
\]
This survival ensures the finer counts for both \(D\) and its
complement in \(A_i\).

Let \(Z\in[0,1]\) be the surviving fraction of a representative
family. If \(\mathbb EZ\ge p\), then
\(\mathbb EZ\le p/2+\Pr[Z\ge p/2]\); hence
\(\Pr[Z\ge p/2]\ge p/2\). Apply this with \(p=u^{-3l}\).
The three subblock partitions are independent, so with further
probability at least \(u^{-24}/8\ge u^{-25}\), all three retain
these fractions. Let \(X_i\) consist of the distinct signed weights of
the surviving smaller masks, retaining their associated masks. For
sufficiently large \(n\),
\begin{equation}\label{surv:representative-count}
 |X_i|\ge \frac{2^{L_i}}{u^{14}}.
\end{equation}

We have now retained many distinct weights whose masks already meet
all cardinality prescriptions. It remains to ensure that the low
modular filters keep compatible choices. The next argument first
bounds collisions under the random moduli, and then exposes the
uniform target residues in the order that fixes the top choice before
its children's choices.

\emph{Hit occupied low-residue classes.}
This collision-to-occupancy argument is the residue-coverage method of
Austrin, Kaski, Koivisto, and Nederlof~\cite[Proposition~3.5]{AKKN16}
and its product-modulus version in Nederlof and
W\k{e}grzycki~\cite[full version, Lemma~3.2]{NW21}; we give the bounds
for the present nested moduli explicitly.
Expose the prime factors forming \(p_i\). For fixed \(X_i\),
the sum of squares of its residue-class sizes modulo \(p_i\) has
expectation at most
\[
 |X_i|+|X_i|^2u^{10}2^{-L_i}
 \le u^{15}|X_i|^2 2^{-L_i},
\]
by \eqref{eq:filter-bounds} and \eqref{surv:representative-count}.
Markov's inequality and a union bound show that, with probability at
least \(1-3/u\), all three sums of squares are at most
\(u^{16}|X_i|^2 2^{-L_i}\). The moduli need not be independent of
one another. Cauchy--Schwarz gives at least
\begin{equation}\label{surv:occupied-residues}
 \frac{2^{L_i}}{u^{16}}\ge \frac{p_i}{u^{16}}
\end{equation}
occupied residues for each \(X_i\).

First expose \(\rho_0\). Assign a surviving smaller mask
\(D_0\subset A_0\) to L and its complement in \(A_0\) to R. As
\(D_0\) varies, the left solution-parent weight is its weight plus
a fixed term: the selected left background weight and the weight of
\(A_L\). Equation~\eqref{surv:occupied-residues} implies that a
uniform \(\rho_0\bmod p_0\) hits one of these weights with
probability at least \(u^{-16}\). On a hit, fix a representative
deterministically, thereby fixing disjoint top masks and parent weights
\(x_L,x_R=t'-x_L\).

Splitting \(A_j\) between the two children cannot change parent \(j\)'s
weight: their union still uses all of \(A_j\). For each side \(j\),
assign a surviving smaller mask
\(D_j\subset A_j\) to its first child and the complement to its
second child. The first-child weight is the weight of \(D_j\) plus
a now fixed term from its background and its portion of the top mask.
Expose the independent residues \(\rho_L,\rho_R\). They both hit
possible first-child residues with probability at least \(u^{-32}\).
The second-child filters then follow automatically from
\(p_j\mid p_0\) and the parent filters. All prescribed leaf counts
hold by the surviving subblock choices; the selected solution has no
background cardinality constraint to satisfy.

The product of the conditional probability lower bounds, for large
\(n\), is at least
\[
 u^{-3}\cdot\frac12\cdot\frac18\cdot u^{-3}
 \cdot u^{-25}\cdot\frac12\cdot u^{-48}
 \ge u^{-100}.
\]
All choices of representatives on hits may be deterministic analysis
choices. They precede both \(q\) and the separating families. For any
subsequent \(q\), one loop iteration has
\(v_L=x_L\bmod P\) and \(v_R=x_R\bmod P\).
\end{proof}
